\section{Pairing with the self-interaction}
\label{sec:gapcontact}

Section~\ref{sec:gap} treated the pairing by the exchange of the Vorton-Scalar alone. This section
adds the self-interaction of the fermion field. It writes out the amplitudes and the gap equation in
full, shows that the exchange of the Vorton-Scalar is itself a self-interaction of finite range,
gives the condition under which the gap equation has a solution that does not depend on the upper
limit of its momentum integral, and gives the gap along the deflation clock when both interactions
act.\SIG{ On the $4+4$ spacetime the contact couplings of form C enter the pair along the hidden axis
with the sign opposite to that of form Q, and the two forms select different pairs.}{}

\subsection{Amplitudes of the four contact channels}

Two particles of the mass $\mst$ with the momenta $\vec p$ and $-\vec p$ in ordinary space scatter
into $\vec p'$ and $-\vec p'$. Let $E=(\mstq+p^{2})^{1/2}$ and $E'$ likewise, and let
$\hh(\vec p)=\sum_j\alpha^{j}p_j+\beta\,\mst$ be the Hamiltonian matrix of one particle, so that
$\tfrac12(1+\hh/E)$ projects on its states of positive energy. The pair along the hidden axis has
the wave function
\begin{equation}
\Phi(\vec p)=\tfrac12\Big(1+\frac{\hh(\vec p)}{E}\Big)\,X^{*}\,\tfrac12\Big(1+\frac{\hh(-\vec p)}{E}\Big)^{T},\qquad X=\Cm\gamma^{0},
\end{equation}
a $16\times16$ matrix. A vertex $\psi^\dagger\,\mathrm{Adj}\,\Theta^{K}\psi$ acting on both particles,
with the adjoint matrix $\mathrm{Adj}$ of the model, gives the amplitude
$\mathrm{Tr}[\Phi(\vec p')^\dagger\,\mathrm{Adj}\,\Theta^{K}\,\Phi(\vec p)\,(\mathrm{Adj}\,\Theta^{K})^{T}]$
divided by the norms of the two wave functions. Summed over the lists $K$ of one rank with the signs
$\sgnK$, this defines the amplitude of the rank $r$. \SIG{In form Q ($\mathrm{Adj}=\beta$, matrices
$\ThQ^{K}$, signs $\sgnQ$), as on the $7+1$ spacetime, the result for the four independent channels of
Sec.~\ref{sec:four} is}{With $\mathrm{Adj}=\beta$ the result for the four independent channels of
Sec.~\ref{sec:four} is}
\begin{equation}
\begin{aligned}
F_0&=\frac{EE'+\mstq-\vec p\cdot\vec p'}{2EE'},\\
F_1&=\frac{3EE'-2\mstq-\vec p\cdot\vec p'}{EE'},\\
F_7&=\frac{3EE'+2\mstq-\vec p\cdot\vec p'}{EE'},\\
F_8&=-\frac{EE'-\mstq-\vec p\cdot\vec p'}{2EE'} .
\end{aligned}
\label{eq:Fclosed}
\end{equation}
$F_0$ is the amplitude $F_S$ of the exchange of the radial mode of the Vorton-Scalar, which couples
to the scalar density, and $F_8=-F_P$ is minus the amplitude of the exchange of its phase mode,
which couples to the pseudoscalar density. At rest $(F_0,F_1,F_7,F_8)=(1,1,5,0)$, the pair coupling
(\ref{eq:Gpair}). The closed forms were verified against the direct evaluation with the 256 matrices
at random momenta, by the program to $1.8\times10^{-15}$ and by the independent script with
projectors built from numerical eigenvectors; \SIG{the $7+1$ spacetime gives the same four
functions.}{form Q of the $4+4$ spacetime gives the same four functions.} They
hold for momenta in ordinary space; a momentum along the hidden axis itself changes them, because
the pair operator points along that axis.
%%<1644>

\emph{The sign of the contact couplings in the two forms.} The same evaluation for the pairs along
$x5,x6,x7$, $\psi^{T}\Cm\gamma^{t}\psi$, and for both forms gives one rule. In form Q the amplitudes of
all four vector pairs along the extra axes are the functions (\ref{eq:Fclosed}). In form C
($\mathrm{Adj}=\Cm$) the amplitudes of the pairs along $x5,x6,x7$ are the same functions, and those of
the pair along the hidden axis are \emph{minus} these functions, at all momenta (direct evaluation
at random momenta, to {{GEN:pc5_dev}}). This is Sec.~\ref{sec:pairs} at general momenta: the pair
along the hidden axis consists of one particle of each half of the field, the pairs along $x5,x6,x7$
of two particles of the same half. The contact amplitudes therefore enter the pairing potential with
a sign
\begin{equation}
\sct=\begin{cases}-1&\text{form C, pair along the hidden axis},\\+1&\text{form C, pairs along }x5,x6,x7,\\+1&\text{form Q, all four pairs}.\end{cases}
\label{eq:sct}
\end{equation}
%%</1644>

The relations (\ref{eq:reduction}) between the nine invariants hold for the amplitudes at rest, and at
general momenta for the part of each amplitude that is even under $\vec p'\to-\vec p'$, which is the
sum of the direct and of the exchange scattering of two identical particles; they do not hold for
the odd part, which does not contribute to a contact interaction in a pair state without orbital
angular momentum. For a strictly local interaction the four couplings (\ref{eq:geff}) therefore
suffice. When the interaction is given a range, as below, the model is \emph{defined} by giving that
range to the four channels.

\subsection{The pairing potential and the gap equation}

The pairing potential of the contact interaction in \SIG{each of these four pair channels}{the
hidden-axis channel} is
\begin{equation}
V_{\mathrm{ct}}(\vec p',\vec p)=\SIG{\sct\,}{}\big(\geff0F_0+\geff1F_1+\geff7F_7+\geff8F_8\big)\,\frac{\mct^{2}}{q^{2}+\mct^{2}},
\label{eq:Vct}
\end{equation}
attractive where it is negative, with $q=|\vec p'-\vec p|$\SIG{ and the sign (\ref{eq:sct})}{}. A contact interaction in a gap equation
needs a regularization at large momentum transfer. The last factor gives each coupling the range
$1/\mct$ of an exchanged particle of the mass $\mct$, the \emph{range parameter}; it is the form that
the exchange of the Vorton-Scalar itself has:
\begin{equation}
V_{\mathrm{VS}}(\vec p',\vec p)=-\gVS^{2}\Big[\frac{F_S}{q^{2}+m_s^{2}}+\frac{F_P}{q^{2}+\mth^{2}}\Big],
\label{eq:VVS}
\end{equation}
with the mass $m_s$ of the radial mode and the mass $\mth$ of the phase mode, which vanishes when the
chiral symmetry is exact\SIG{; it is the same in all four pair channels and in both forms, because the
Yukawa coupling contains the adjoint $\bar\psi$, in which the two halves of the field enter alike}{}.
Comparing Eqs.~(\ref{eq:Vct}) and (\ref{eq:VVS}): \emph{the exchange of the
Vorton-Scalar is a self-interaction of finite range}\SIG{ of form Q}{}, with the momentum-dependent couplings
\begin{equation}
\geff0(q)=-\frac{\gVS^{2}}{q^{2}+m_s^{2}},\qquad\geff8(q)=+\frac{\gVS^{2}}{q^{2}+\mth^{2}} .
\label{eq:gq_exchange}
\end{equation}
At momenta far below $m_s$ the radial mode acts as the scalar contact coupling
$g_0=-\gVS^{2}/m_s^{2}$\SIG{ of form Q}{}. In the phase with unbroken symmetry the two masses are equal and
$\geff0(q)+\geff8(q)=0$, the condition (\ref{eq:chiralcond}) of chiral invariance\SIG{ in form Q}{}.
\SIG{In the $4+4$ model of this paper the contact terms are of form C and the exchange is not: in the
pair along the hidden axis the two enter with opposite relative sign, and in the pairs along
$x5,x6,x7$ with the same. }{}The program confirms
that the scalar contact coupling $-\gVS^{2}/m_s^{2}$ with $\mct=m_s$ reproduces the kernel of the
exchange of the radial mode to $3\times10^{-16}$.

The gap function $\Delta(k)$ of the paired state obeys, at the temperature $T$,
\begin{equation}
\begin{aligned}
\Delta(k)&=\frac{1}{4\pi^{2}}\int_0^{\kmax}dk'\,k'^{2}\,W(k,k')\,\frac{\Delta(k')}{E_{k'}}\tanh\frac{E_{k'}}{2T},\\
W(k,k')&=\frac12\int_{-1}^{1}d\cos\theta\,\big[-V_{\mathrm{VS}}-V_{\mathrm{ct}}\big],
\end{aligned}
\label{eq:gapeq}
\end{equation}
with $\xi_k=(k^{2}+m^{2})^{1/2}-\mu$, the quasi-particle energy $E_k=(\xi_k^{2}+\Delta(k)^{2})^{1/2}$,
the chemical potential $\mu=(\kF^{2}+m^{2})^{1/2}$, and the angle $\theta$ between $\vec p$ and
$\vec p'$ ($k=|\vec p|$, $k'=|\vec p'|$); $\Delta_F=\Delta(\kF)$; $\kmax$ is the upper limit of the
momentum integral, which the next subsection examines. It is the condition that the CJT
effective action be stationary with respect to the pair part of the propagator, in the truncation
of Sec.~\ref{sec:cjt}, and equally the form that the Kohn-Sham equations take with a pairing field.
The angular integral is elementary, because each amplitude is linear in $\cos\theta$.

\subsection{The upper limit of the momentum integral: an ultraviolet criterion}
\label{sec:uv}

At momenta far above all masses every term of the pairing potential has the form of a coupling
times a function of $\cos\theta$ divided by $q^{2}$, with $F_0,F_S,F_P\to\tfrac12(1-\cos\theta)$,
$F_8\to-\tfrac12(1-\cos\theta)$ and $F_1,F_7\to3-\cos\theta$. The kernel is then homogeneous,
$W(k,k')\to w_\infty(k'/k)/k^{2}$ with a function $w_\infty$ of the ratio of the two momenta, and
above the gap Eq.~(\ref{eq:gapeq}) at zero temperature reads
$\Delta(k)=\tfrac1{4\pi^{2}}\int_0^\infty dt\,t\,w_\infty(t)\,\Delta(tk)$. A power
$\Delta(k)\propto k^{-s_\Delta}$ solves it if the exponent $s_\Delta$ obeys
\begin{equation}
\int_0^\infty dt\;t^{1-s_\Delta}\,w_\infty(t)=4\pi^{2} .
\label{eq:uvroot}
\end{equation}
Because $W(k,k')$ is symmetric in its arguments, $w_\infty(1/t)=t^{2}\,w_\infty(t)$; the left-hand
side of Eq.~(\ref{eq:uvroot}) is therefore unchanged when $s_\Delta$ is replaced by $2-s_\Delta$, and
it is smallest at $s_\Delta=1$. Its value there is the \emph{ultraviolet strength} of the kernel,
\begin{equation}
\Uuv=\int_0^\infty dt\;w_\infty(t).
\label{eq:Uuvdef}
\end{equation}
If $\Uuv<4\pi^{2}$, Eq.~(\ref{eq:uvroot}) has two real roots, $s_\Delta>1$ and $2-s_\Delta<1$; the gap
function falls as $k^{-s_\Delta}$ with the larger root, and the gap equation has a solution that
does not depend on $\kmax$, which is approached as $\kmax^{-2(s_\Delta-1)}$, slowly when $\Uuv$ is
close to $4\pi^{2}$. If $\Uuv>4\pi^{2}$ the roots are complex, no such solution exists, and with an
upper limit the gap grows in proportion to $\kmax$: the pair collapses to short distances. This is
the fall to the center of relativistic particles in an attractive Coulomb potential whose strength
exceeds a critical value; the exchange of a boson of the mass $m_s$ gives a Yukawa potential, which
is a Coulomb potential at distances below $1/m_s$.

The integrals over $t$ and over the angle are elementary. With
$\int_0^\infty dt/(1+t^{2}-2t\cos\theta)=(\pi-\theta)/\sin\theta$ the two functions of the angle
give
\begin{equation}
\begin{aligned}
u_S&\equiv\int_0^\infty\!dt\,\frac12\int_{-1}^{1}\!d\cos\theta\,\frac{\tfrac12(1-\cos\theta)}{1+t^{2}-2t\cos\theta}=\frac{\pi^{2}}{8}-\frac12,\\
u_V&\equiv\int_0^\infty\!dt\,\frac12\int_{-1}^{1}\!d\cos\theta\,\frac{3-\cos\theta}{1+t^{2}-2t\cos\theta}=\frac{3\pi^{2}}{4}-1,
\end{aligned}
\label{eq:uSuV}
\end{equation}
that is $u_S={{GENM:uv_IS}}$ and $u_V={{GENM:uv_IV}}$, and therefore
\begin{equation}
\begin{aligned}
\Uuv={}&2\,u_S\,\gVS^{2}-\SIG{\sct\,}{}\mct^{2}\Big[\big(\geff0-\geff8\big)\,u_S\\
&+\big(\geff1+\geff7\big)\,u_V\Big]
\end{aligned}
\label{eq:Uuv}
\end{equation}
when both modes of the Vorton-Scalar are exchanged (half the first term for the radial mode alone).
The formula was verified against the numerical integral of the kernel by the program and by the
independent script. \emph{The exchange of the two modes alone reaches $\Uuv=4\pi^{2}$ at
$\gVS^{2}=4\pi^{2}/(2u_S)={{GENM:uv_g2c_exchange}}$}, the clock reading
$\td={{GENM:uv_tau_exchange}}$ of the histories of Sec.~\ref{sec:gap}, beyond their last reading
{{GEN:uv_tau_last}}. A scalar contact attraction \SIG{with $\sct=+1$ }{}brings the criterion nearer, a
vector repulsion moves it away (Table~\ref{tab:uv})\SIG{; with $\sct=-1$ it is the other way round, and
the vector contact coupling $+0.3\,\gVS^{2}/m^{2}$ in the hidden-axis pair of form C reaches the
criterion at $\gVS^{2}={{GENM:uv_g2c_v3minus}}$, the clock reading ${{GENM:uv_tau_v3minus}}$, from which
on no gap is given}{}.

\begin{table*}
\caption{\label{tab:uv}The ultraviolet strength $\Uuv$ of the pairing kernel per unit $\gVS^{2}$, the
value of $\gVS^{2}$ at which it reaches $4\pi^{2}$, and the clock reading of that moment
($\gVS=e^{\td/2}$), for the exchange of the two modes of the Vorton-Scalar alone and with contact
couplings ($\mct=m$)\SIG{, for the two signs (\ref{eq:sct})}{}. A dash: the criterion is never reached.}
\begin{ruledtabular}
\begin{tabular}{\SIG{>{\raggedright\arraybackslash}p{0.3\textwidth}cccccc}{lccc}}
interaction & \SIG{$\sct=+1$: $\Uuv/\gVS^{2}$}{$\Uuv/\gVS^{2}$} & critical $\gVS^{2}$ & $\td$\SIG{ & $\sct=-1$: $\Uuv/\gVS^{2}$ & critical $\gVS^{2}$ & $\td$}{} \\
\hline
{{ROWS:tbl_fm_uv\SIG{}{|cols=0,1,2,3}}}
\end{tabular}
\end{ruledtabular}
\end{table*}

The criterion was tested numerically: the gap was computed with the upper limits $\kmax=60$,
$2\times10^{3}$, $2\times10^{4}$ and $2\times10^{5}$, in units of the larger of $m$ and $\kF$, for
{{GEN:uv_n_cases}} kernels (four sets of $m_s$ and $\kF$, four clock readings, seven combinations of
couplings, and both signs of the contact part\SIG{}{; the sign opposite to that of this spacetime
occurs on the $4+4$ spacetime}). Below the criterion the gap converges
(Table~\ref{tab:uvclasses}); above it ({{GEN:uv_n_above}} kernels\SIG{}{, all with the opposite sign})
it grows by a factor between {{GEN:uv_growth_lo}} and {{GEN:uv_growth_hi}} from $\kmax=2\times10^{4}$
to $2\times10^{5}$, so that its value is not a property of the model. \emph{All gaps of this paper are
computed with $\kmax=2\times10^{4}$, and none is given above the criterion.} A limit of 60, with
which gaps of this model were computed before, gives gaps of the exchange alone that are too small
by up to {{GEN:uv_exchange_error}} per cent at the strongest coupling of the histories, and gaps with
contact couplings that are too small by the percentages of the last column of the table.

\begin{table}
\caption{\label{tab:uvclasses}Convergence of the gap with the upper limit $\kmax$ of the momentum
integral, by classes of the ultraviolet strength (kernels below the criterion with a gap above
$10^{-4}$ of the Fermi kinetic energy): the largest relative change of the gap from
$\kmax=2\times10^{4}$ to $2\times10^{5}$, and the largest error of the limit $\kmax=60$.}
\begin{ruledtabular}
\begin{tabular}{lccc}
$\Uuv/4\pi^{2}$ & kernels & change & error of 60 (per cent) \\
\hline
{{ROWS:tbl_fm_uv_classes}}
\end{tabular}
\end{ruledtabular}
\end{table}

\subsection{The gap along the clock with contact couplings}

All couplings of the four-dimensional description grow as $e^{\td}$ (Sec.~\ref{sec:clock}), so that
the ratio of a contact coupling to $\gVS^{2}$ is fixed. Table~\ref{tab:gapcontact} and
Fig.~\ref{fig:paircontact} give the gap\SIG{, for the sign $\sct=+1$,}{} for seven choices: the exchange alone; the exchange with a
scalar contact attraction $\geff0=-\gVS^{2}/m^{2}$; with a scalar attraction and a vector repulsion
one tenth as large, $\geff1=+0.1\,\gVS^{2}/m^{2}$; with a vector repulsion alone,
$\geff1=+0.3\,\gVS^{2}/m^{2}$; with an axial coupling $\geff7=-0.1\,\gVS^{2}/m^{2}$; the contact
couplings without the Vorton-Scalar; and the exchange alone with a phase mode of the mass
$\mth=0.3\,m$. The range parameter is $\mct=m$.

\begin{table*}
\caption{\label{tab:gapcontact}The gap at the Fermi surface in units of the Fermi kinetic energy,
with contact couplings\SIG{ that enter with the sign $\sct=+1$ (form Q, all four vector pairs along the
extra axes; form C, the pairs along $x5,x6,x7$)}{}, at the clock reading {{GEN:gapc_tau}}
[$\gVS={{GENM:gapc_g}}$], for the Fermi
momentum $\kF=0.5\,m$ and three masses of the radial mode. Contact couplings in units of
$\gVS^{2}/m^{2}$; $\mct=m$.}
\begin{ruledtabular}
\begin{tabular}{lccc}
interaction & $m_s=0.1\,m$ & $m_s=0.3\,m$ & $m_s=m$ \\
\hline
{{ROWS:tbl_gap_contact}}
\end{tabular}
\end{ruledtabular}
\end{table*}

\begin{figure*}
\includegraphics[width=\textwidth]{fig_pair_contact}
\caption{\label{fig:paircontact}The gap \SIG{with contact couplings that enter with the sign $\sct=+1$
(form Q; form C in the pairs along $x5,x6,x7$)}{of the hidden-axis pair} along the deflation clock\SIG{}{ with contact
couplings}, for the Fermi momentum $0.5\,m$ and three masses $m_s$ of the radial mode of the
Vorton-Scalar. Upper row: the exchange of the Vorton-Scalar alone; with a scalar contact attraction;
with a scalar attraction and a vector repulsion one tenth as large; the contact couplings alone.
Lower row: the exchange alone; with a vector contact repulsion; with an axial contact coupling; with
a phase mode of the mass $0.3\,m$. Couplings and range parameter as in
Table~\ref{tab:gapcontact}.}
\end{figure*}

Three statements summarize the result\SIG{ for $\sct=+1$}{}; each was verified at every clock reading at
which the gap exceeds $10^{-6}$ of the Fermi kinetic energy.

\emph{A scalar contact attraction raises the gap}, strongly so where the exchange alone is weak: for
$m_s=m$ at the clock reading of Table~\ref{tab:gapcontact} by the factor {{GEN:gapc_factor}}.

\emph{A vector contact repulsion lowers it.} A vector coupling of three tenths of $\gVS^{2}/m^{2}$
removes most of the gap of a weakly paired state (Table~\ref{tab:gapcontact}, $m_s=m$) and reduces
the gap of a strongly paired one.

\emph{A mass of the phase mode lowers it}, because it weakens the attraction of that mode at small
momentum transfer.

Without contact couplings the program reproduces the gaps of Sec.~\ref{sec:gap} to $7\times10^{-16}$.
For small gaps the critical temperature remains $0.567$ times the gap. An independent solution of
Eq.~(\ref{eq:gapeq}) on another grid of momenta, by plain iteration, gives the gaps of
Table~\ref{tab:gapcontact} for the first three rows to {{GEN:gapc_indep}} (relative). The law of growth
(\ref{eq:growth}) holds with contact couplings as it does without, because every part of the kernel
grows as $e^{\td}$.
%%<1644>

\subsection{The pair along the hidden axis in form C}
\label{sec:gapformc}

In the definition of the model, form C, the contact couplings enter the pair along the hidden axis
with $\sct=-1$. Table~\ref{tab:gapsign} and Fig.~\ref{fig:gapsign} compare the two signs for the same
couplings. With $\sct=-1$ the couplings act in the opposite way: \emph{a scalar contact attraction
lowers the gap} (for $m_s=m$ it removes it at the clock reading of the table), an axial coupling
$\geff7<0$ lowers it, and \emph{a vector contact repulsion raises it}; the contact couplings of the
signs of set B alone give no gap in this pair at any clock reading. The gap with $\sct=-1$ and the
vector coupling agrees with an independent solution of the gap equation to {{GEN:gapc_indep_minus}}.

\begin{table*}
\caption{\label{tab:gapsign}The gap at the Fermi surface in units of the Fermi kinetic energy at the
clock reading {{GEN:fm_gap_tau}}, for $\kF=0.5\,m$ and two masses of the radial mode, for the two
signs with which the contact couplings enter the pairing potential: $\sct=+1$ (form Q; form C in the
pairs along $x5,x6,x7$) and $\sct=-1$ (form C in the pair along the hidden axis). Last two columns:
the ultraviolet strength of the kernel for $m_s=m$. Contact couplings in units of $\gVS^{2}/m^{2}$;
$\mct=m$.}
\begin{ruledtabular}
\begin{tabular}{>{\raggedright\arraybackslash}p{0.3\textwidth}cccccc}
interaction & $m_s=0.3\,m$, $\sct=+1$ & $-1$ & $m_s=m$, $\sct=+1$ & $-1$ & $\Uuv/4\pi^{2}$, $\sct=+1$ & $-1$ \\
\hline
{{ROWS:tbl_fm_gap}}
\end{tabular}
\end{ruledtabular}
\end{table*}

\begin{figure*}
\includegraphics[width=0.86\textwidth]{fig_forms_gap}
\caption{\label{fig:gapsign}The gap at the Fermi surface, in units of the Fermi kinetic energy, along
the deflation clock, for $m_s=m$ and $\kF=0.5\,m$. Left: with the scalar contact coupling
$\geff0=-\gVS^{2}/m^{2}$; right: with the vector contact coupling $\geff1=+0.3\,\gVS^{2}/m^{2}$. Dotted
line: the exchange of the Vorton-Scalar alone. Dashed line (squares): the contact coupling with the
sign $\sct=+1$ (form Q; form C in the pairs along $x5,x6,x7$). Solid line (circles): with the sign
$\sct=-1$ (form C in the pair along the hidden axis). In the shaded band the kernel with $\sct=-1$
exceeds the ultraviolet criterion of Sec.~\ref{sec:uv}, and no gap is given.}
\end{figure*}

\emph{In form C the self-interaction removes the degeneracy of the pair along the hidden axis with
the pairs along $x5,x6,x7$, and so selects the direction of the order parameter}
(Table~\ref{tab:gapwins}). With an attractive scalar coupling the pairs along the additional
time-like axes have the larger gap, not the pair along the hidden axis; with a repulsive vector
coupling alone the pair along the hidden axis has the larger gap. In form Q, as on the $7+1$
spacetime, the four pairs have the same gap, and the direction is free (Sec.~\ref{sec:order}).

\begin{table}
\caption{\label{tab:gapwins}Form C: the number of clock readings, among those at which both gaps are
defined, at which the pairs along $x5,x6,x7$ ($\sct=+1$) or the pair along the hidden axis
($\sct=-1$) have the larger gap (four sets of $m_s$ and $\kF$).}
\begin{ruledtabular}
\begin{tabular}{p{0.52\columnwidth}cc}
interaction & $x5,x6,x7$ & hidden axis \\
\hline
{{ROWS:tbl_fm_gap_wins}}
\end{tabular}
\end{ruledtabular}
\end{table}
%%</1644>
